\chapter{Is the difference significant}
\label{ch:statistics}

\section{What was run}

A paired sign flip permutation clustered at the template level, exhaustive over
the full arrangement space, with step up correction across a family of
\vFamilySize{} comparisons, and an absolute practical effect gate. The design
quantities are in Table~\ref{tab:design}.

The test split holds \vTestTemplates{} templates, which gives \vArrangements{}
arrangements and a smallest attainable two sided p value of \vMinAttainableP{},
comfortably below the pre-registered level of \vAlpha{}. \textbf{Every
comparison in this analysis is decidable}, which was not true of the first
measurement this project made and is the whole point of the corpus repair
described in Chapter~\ref{ch:corpus}.

\section{The outcome categories, including the empty ones}

\begin{table}[htbp]
  \centering
  \caption{Every verdict category in the family of comparisons, including the
  two that came back empty. A summary listing only the categories that occurred
  would drop them on exactly the runs where nothing landed in them.}
  \label{tab:outcomes}
  \begingroup\small\setlength{\tabcolsep}{4pt}
\begin{tabular}{lr}
\toprule
\textbf{Verdict} & \textbf{Comparisons} \\
\midrule
improvement & 11 \\  % results: experiments/results/analysis/2026-08-26T21-20-11Z_p6-analysis_6457ac7/analysis.json
regression & 5 \\  % results: experiments/results/analysis/2026-08-26T21-20-11Z_p6-analysis_6457ac7/analysis.json
significant but below the practical threshold & 0 \\  % results: experiments/results/analysis/2026-08-26T21-20-11Z_p6-analysis_6457ac7/analysis.json
no significant change & 17 \\  % results: experiments/results/analysis/2026-08-26T21-20-11Z_p6-analysis_6457ac7/analysis.json
inconclusive: the design cannot reach alpha & 0 \\  % results: experiments/results/analysis/2026-08-26T21-20-11Z_p6-analysis_6457ac7/analysis.json
\midrule
total & 33 \\  % results: experiments/results/analysis/2026-08-26T21-20-11Z_p6-analysis_6457ac7/analysis.json
\bottomrule
\end{tabular}
\endgroup
\par\smallskip\noindent{\footnotesize Source: \rpath{experiments/results/analysis/2026-08-26T21-20-11Z_p6-analysis_6457ac7/analysis.json}}

\end{table}

Table~\ref{tab:outcomes} carries all five categories. Two are worth pointing at
because they are empty.

\textbf{Significant but below the practical threshold: \vBelowThreshold{}.} This
category exists because a difference can be certain and still be too small to
act on, and a reporting scheme that silently merges it into either neighbour is
a reporting scheme that will one day sell a rounding error as a result. On this
family nothing landed in it. It is printed anyway.

\textbf{Inconclusive because the design cannot reach alpha: \vUnderpowered{}.}
This is the category that held all \vOldFamilySize{} comparisons of the first
analysis, when the test split had \vOldClusters{} clusters and a smallest
attainable p of \vOldMinAttainableP{}. It is now empty, which is the measurable
result of the repair.

\section{What is certified, and what is not}

\begin{table}[htbp]
  \centering
  \caption{All comparisons in the family, read in the direction they were
  measured. The full table is Table~\ref{tab:comparisons} in
  Appendix~\ref{app:tables}; this one repeats its shape for reference.}
  \label{tab:cross-check}
  \begingroup\small\setlength{\tabcolsep}{4pt}
\begin{tabular}{lr}
\toprule
\textbf{Cross check} & \textbf{Comparisons} \\
\midrule
Comparisons checked & 33 \\  % results: experiments/results/analysis/2026-08-26T21-20-11Z_p6-analysis_6457ac7/analysis.json
Verdicts that agree & 33 \\  % results: experiments/results/analysis/2026-08-26T21-20-11Z_p6-analysis_6457ac7/analysis.json
Verdicts that disagree & 0 \\  % results: experiments/results/analysis/2026-08-26T21-20-11Z_p6-analysis_6457ac7/analysis.json
\bottomrule
\end{tabular}
\endgroup
\par\smallskip\noindent{\footnotesize Source: \rpath{experiments/results/analysis/2026-08-26T21-20-11Z_p6-analysis_6457ac7/analysis.json}}

\end{table}

\textbf{Certified against the language model, after correction:} the rule table
leads it overall and on the seen locales, and both classical engines lead it on
clean and on partial markup. The language model's missing declaration recall is
certified better than both, and its precision on both accusation codes is
certified worse than both.

\textbf{Not certified:} everything on the hostile tier, everything on the held
out locale, and the language model against the n-gram model overall. Those are
precisely the slices where the language model was expected to do well, and the
honest report of them is that the differences are not large enough to call at
\vClusters{} clusters, in either direction.

\textbf{Not the same statement as no difference.} Comparisons in the \emph{no
significant change} category are comparisons where this design, with this cluster
count, could not resolve the effect. Reading them as evidence of equivalence
would be reading the absence of a certificate as a certificate.

The complete list, every comparison with its effect, its exact p value, its
adjusted p value and its verdict, is Table~\ref{tab:comparisons}.

\section{One reading hazard, stated explicitly}

The comparison names in the result file read \code{baseline\_vs\_candidate} and
the effect is computed as candidate minus baseline. A comparison named for the
language model against the rule table with a verdict of \emph{improvement}
therefore means \textbf{the rule table improved over the language model}, not the
reverse.

The file is self describing: every comparison carries its baseline and its
candidate by name. But the string reads backwards at a glance, which is exactly
the kind of thing that ends up in a report as a sentence pointing the wrong way.
The table generator never renders the comparison name as a sentence. It reads the
baseline, the candidate and the effect and builds the sentence, and the
\emph{Comparison} column of Table~\ref{tab:comparisons} is that sentence.

\section{The family size was fixed in advance, and it cost something}

Every comparison between the two classical engines carries a larger adjusted p
value in this analysis than the same comparison carried when it was corrected
inside the smaller family of the previous phase. The underlying p values did not
move. The earlier analysis file is untouched and is still in the repository.

That was registered before the runs, with the reason: the comparison this phase
existed to make was the one against the language model, and a family containing
only that pair would have been a third the size and would have produced smaller
adjusted p values for exactly the comparisons the project most wanted to be able
to report. Fixing the family size while it was still costless is the only way
that sentence can be written afterwards without it being an excuse.

\section{The harness's own verdicts, as a cross check}

The external harness carries its own end to end verdict machinery. Its
statistical gate and its underpowered gate are the ones used above. Its practical
gate reads a relative percentage where this project pre-registered an absolute
difference, so a disagreement between the two verdict columns would be that gap
and nothing else.

It was run alongside on the same numbers as a visible cross check.
\vCrossCheckChecked{} comparisons were checked and \vCrossCheckDisagreements{}
disagreed. That gap therefore changed no verdict on this data, which is worth
saying plainly and is not a reason to withdraw the filed issue: a gap that happens
not to bite on one dataset is still a gap, and the dataset it would bite on is any
slice with a small baseline.

\section{The secondary analysis, and why nothing cites it}

A secondary analysis clustered by form rather than by template reaches
significance easily. It is computed, it is in the result file, and it is labelled
anticonservative everywhere it appears, because clustering by form asserts that
two locales of one template are independent, which is a stronger claim than this
project's design makes.

Nothing in this report, in the README, or in any document in this repository
cites it as evidence. It is present so that the difference between the two
clusterings is visible rather than assertable, and that is the only reason it is
present.
